\documentclass[10pt, a4paper]{letter}
\usepackage{ngerman}
\usepackage[utf8]{inputenc}
\usepackage{nopageno}
\usepackage{xcolor}
\usepackage[official]{eurosym}

% Tabellen
\usepackage{tabularx}
\setlength{\extrarowheight}{5pt}

% Montserrat als Schriftart (i like it)
\usepackage[defaultfam,tabular,lining]{montserrat}
\usepackage[T1]{fontenc}
\renewcommand*\oldstylenums[1]{{\fontfamily{Montserrat-TOsF}\selectfont #1}}

% Seitenränder und Alignment
\usepackage[left=1.5cm, right=1.5cm, top=2cm, bottom=2cm]{geometry}
\usepackage[document]{ragged2e}

% Metadata
\usepackage[pdftex,
            pdfauthor={Christopher Katins},
            pdftitle={Mieterselbstauskunft}]{hyperref}


% Normale Hyperref Text Fields können nicht ohne Weiteres mit /hspace o. Ä. auf die maximale Linewidth gezogen werden, daher diese super hilfreiche Neudefinition, die genau das macht!
% Quelle: https://tex.stackexchange.com/questions/77629/defining-a-width-that-fills-remaining-horizontal-space-for-text-fields-in-hyperr
\newlength\TextFieldLength
\newcommand\TextFieldFill[2][]{%
  \setlength\TextFieldLength{\linewidth}%
  \settowidth{\dimen0}{#2 }%
  \addtolength\TextFieldLength{-\dimen0}%
  \addtolength\TextFieldLength{-2.22221pt}%
  \TextField[#1,width=\TextFieldLength]{\raisebox{2pt}{#2 }}%
}
% To copy: \TextFieldFill[name=]{}


\begin{document}
\justifying

% Globales Design der TextFields
\def\DefaultHeightofText{12pt}
\def\DefaultOptionsofText{backgroundcolor=gray!30,borderwidth=0,bordercolor=white}
\def\DefaultOptionsofCheckBox{backgroundcolor=gray!30,borderwidth=0,bordercolor=white}

\textbf{\huge Selbstschuldnerische Mietbürgschaft}
\vspace{.4cm}

\textbf{\Large Zum Mietobjekt}

\TextFieldFill[name=str-nmr]{Straße, Nr. (Zimmer, Etage, ...): }

\TextFieldFill[name=plz-ort]{PLZ, Ort: }

\TextFieldFill[name=mietbeginn]{Mietbeginn: }

\TextFieldFill[name=nettokaltmiete]{Nettokaltmiete: }

\vspace{.8cm}
\textbf{\Large Zum Vermieter}

\TextFieldFill[name=name]{Name, Vorname: }

\TextFieldFill[name=str-nmr-bisherig]{Straße, Nr. (bisherig): }

\TextFieldFill[name=plz-ort-bisherig]{PLZ, Ort (bisherig): }

\vspace{.8cm}
\textbf{\Large Zum ersten Hauptmieter}

\TextFieldFill[name=name]{Name, Vorname: }

\TextFieldFill[name=str-nmr-bisherig]{Straße, Nr. (bisherig): }

\TextFieldFill[name=plz-ort-bisherig]{PLZ, Ort (bisherig): }

\vspace{.8cm}
\textbf{\Large Zum zweiten Hauptmieter}

\TextFieldFill[name=name]{Name, Vorname: }

\TextFieldFill[name=str-nmr-bisherig]{Straße, Nr. (bisherig): }

\TextFieldFill[name=plz-ort-bisherig]{PLZ, Ort (bisherig): }

\vspace{.8cm}
\textbf{\Large Zum ersten Bürge}

\TextFieldFill[name=name]{Name, Vorname: }

\TextFieldFill[name=str-nmr-bisherig]{Straße, Nr. (bisherig): }

\TextFieldFill[name=plz-ort-bisherig]{PLZ, Ort (bisherig): }

\TextFieldFill[name=geb-datum]{Geburtsdatum: }

\TextFieldFill[name=perso-nummer]{Personalausweisnummer: }

\vspace{.8cm}
\textbf{\Large Zum zweiten Bürge}

\TextFieldFill[name=name]{Name, Vorname: }

\TextFieldFill[name=str-nmr-bisherig]{Straße, Nr. (bisherig): }

\TextFieldFill[name=plz-ort-bisherig]{PLZ, Ort (bisherig): }

\TextFieldFill[name=geb-datum]{Geburtsdatum: }

\TextFieldFill[name=perso-nummer]{Personalausweisnummer: }



\vspace{.8cm}
\textbf{\Large Umfang der Bürgschaft}

Hiermit übernehmen wir, die oben genannten Bürgen, gegenüber dem Vermieter die \textbf{selbstschuldnerische Bürgschaft}
für alle bestehenden und zukünftigen Verbindlichkeiten des Mieters aus dem oben genannten Mietverhältnis als
\textbf{Gesamtschuldner} (§ 421 BGB). Das bedeutet, der Vermieter kann jeden Bürgen einzeln für die gesamte Summe in
Anspruch nehmen. Die Bürgschaft umfasst insbesondere:
\begin{itemize}
    \item Mietzahlungen (Kaltmiete und Nebenkostenvorauszahlungen).
    \item Nachzahlungen aus Betriebskostenabrechnungen.
    \item Schäden an der Mietsache und unterlassene Schönheitsreparaturen.
    \item Sämtliche sonstigen Forderungen aus dem Mietvertrag (z. B. Verzugskosten, Räumungskosten).
\end{itemize}
Wir verzichten hiermit ausdrücklich auf die \textbf{Einrede der Vorausklage} (§ 771 BGB) sowie auf die Einrede der
Anfechtbarkeit und Aufrechenbarkeit (§ 770 BGB). Wir verpflichten uns, auf erste Anforderung des Vermieters zu leisten.
Diese Bürgschaft ist unbefristet und gilt auch bei einer Erhöhung der Miete sowie bei einem Vermieterwechsel.

\textbf{Haftungsbegrenzung:} Die Bürgschaft ist der Höhe nach begrenzt auf einen Höchstbetrag von: \\
\TextField[name=unterschrift,width=3cm]{} Euro. (Hinweis: Max. drei Nettokaltmieten gemäß § 551 BGB)

\vspace{3em}

\begin{minipage}[t]{8.5cm}
    \TextField[name=ort_datum_mieter,width=8.5cm]{}\\
    Ort, Datum
    \vspace{3em} % Platz für die eigentliche Unterschrift

    \rule{8.5cm}{0.5pt} \\
    Unterschrift erster Hauptmieter
    \vspace{3em} % Platz für die eigentliche Unterschrift

    \rule{8.5cm}{0.5pt} \\
    Unterschrift zweiter Hauptmieter
    \vspace{3em} % Platz für die eigentliche Unterschrift

    \rule{8.5cm}{0.5pt} \\
    Unterschrift erster Bürge
    \vspace{3em} % Platz für die eigentliche Unterschrift

    \rule{8.5cm}{0.5pt} \\
    Unterschrift zweiter Bürge
\end{minipage}%
\hfill
\begin{minipage}[t]{8.5cm}
    \TextField[name=ort_datum_vermieter,width=8.5cm]{}\\
    Ort, Datum
    \vspace{3em} % Platz für die eigentliche Unterschrift

    \rule{8.5cm}{0.5pt} \\
    Unterschrift Vermieter
\end{minipage}

\end{document}
